% ============================================================
%  Informe de actividad Fase 2 (Módulo 2) — Técnicas Avanzadas de IA
%  Universidad Autónoma de Occidente
%  Formato IEEE Conference
% ============================================================

\documentclass[conference]{IEEEtran}
\IEEEoverridecommandlockouts

% --- Codificación y fuentes ---
\usepackage[utf8]{inputenc}
\usepackage[T1]{fontenc}
\usepackage[spanish,es-noshorthands]{babel}

% --- Matemáticas ---
\usepackage{amsmath,amssymb,amsfonts}

% --- Gráficos y diagramas ---
\usepackage{graphicx}
\usepackage{tikz}
\usetikzlibrary{shapes.geometric, arrows.meta, positioning, backgrounds, fit, calc}

% --- Tablas ---
\usepackage{booktabs}
\usepackage{tabularx}
\usepackage{multirow}
\usepackage{array}

% --- Código fuente ---
\usepackage{listings}
\usepackage{xcolor}

% --- Tipografía y formato ---
\usepackage{microtype}
\usepackage{float}
\usepackage{cite}
\usepackage{url}
\usepackage{hyperref}

% ============================================================
%  Colores y estilos de listings
% ============================================================
\definecolor{clrGreen}{rgb}{0,0.55,0}
\definecolor{clrGray}{rgb}{0.5,0.5,0.5}
\definecolor{clrPurple}{rgb}{0.50,0,0.70}
\definecolor{clrBg}{rgb}{0.97,0.97,0.97}
\definecolor{clrOrange}{rgb}{0.80,0.40,0}
\definecolor{clrBlue}{rgb}{0.10,0.20,0.70}

\lstdefinelanguage{yaml}{
    keywords={true,false,null,yes,no},
    keywordstyle=\color{clrBlue}\bfseries,
    sensitive=false,
    comment=[l]{\#},
    commentstyle=\color{clrGreen}\itshape,
    morestring=[b]',
    morestring=[b]",
    stringstyle=\color{clrOrange},
}

\lstdefinelanguage{json}{
    morestring=[b]",
    stringstyle=\color{clrOrange},
    literate=
        *{0}{{{\color{clrPurple}0}}}{1}
         {1}{{{\color{clrPurple}1}}}{1}
         {2}{{{\color{clrPurple}2}}}{1}
         {3}{{{\color{clrPurple}3}}}{1}
         {4}{{{\color{clrPurple}4}}}{1}
         {5}{{{\color{clrPurple}5}}}{1}
         {:}{{{\color{black}{:}}}}{1}
         {,}{{{\color{black}{,}}}}{1},
}

\lstdefinestyle{estiloBase}{
    backgroundcolor=\color{clrBg},
    commentstyle=\color{clrGreen},
    keywordstyle=\color{clrBlue}\bfseries,
    numberstyle=\tiny\color{clrGray},
    stringstyle=\color{clrPurple},
    basicstyle=\footnotesize\ttfamily,
    breakatwhitespace=false,
    breaklines=true,
    captionpos=b,
    keepspaces=true,
    numbers=left,
    numbersep=4pt,
    showspaces=false,
    showstringspaces=false,
    tabsize=2,
    frame=single,
    rulecolor=\color{black!30},
}
\lstset{style=estiloBase}

\newcommand{\code}[1]{\texttt{\small #1}}

\hypersetup{
    colorlinks=true,
    linkcolor=blue!70!black,
    citecolor=blue!70!black,
    urlcolor=blue!70!black,
}

% ============================================================
%  DOCUMENTO
% ============================================================
\begin{document}

\title{Agente Conversacional con Web Scraping, Recuperación Densa\\
y Memoria Persistente para Información Institucional Hospitalaria:\\
Fundación Valle del Lili}

\author{
    \IEEEauthorblockN{%
        Frank Edward Daza González\quad
        Andrés Fernando López Rendón\quad
        Viviana Valle Fernández%
    }
    \IEEEauthorblockA{%
        Técnicas Avanzadas de IA Aplicadas en Modelos de Lenguaje\\
        Universidad Autónoma de Occidente, Cali, Colombia\\
        Docente: Jan Polanco Velasco
    }
}

\maketitle

% ============================================================
\begin{abstract}
% ============================================================
Este artículo presenta la segunda fase del sistema de asistencia conversacional
para la Fundación Valle del Lili, que evoluciona del pipeline RAG léxico (BM25)
del Módulo~1 hacia un agente conversacional completo con tres componentes
nucleares: recuperación densa en Qdrant con \textit{embeddings}, memoria
multi-turno en PostgreSQL y orquestación dinámica con LangGraph. El
\textit{web scraping} del corpus institucional, descrito en la fase anterior,
mantiene su función como fuente de verdad textual y alimenta el proceso de
ingesta hacia el almacén vectorial. El agente orquesta tres herramientas
complementarias: \code{faq\_estructurada} (respuestas deterministas sobre un
JSON curado), \code{rag\_denso} (recuperación semántica por similitud coseno
con diversificación MMR y reranking cruzado opcional) y
\code{listar\_estructurado} (agregación y conteo por filtros de \textit{payload}
en Qdrant). La transmisión de la respuesta al cliente se realiza mediante
eventos SSE extendidos (\textit{pensamiento}, \textit{herramienta},
\textit{token}, \textit{fuentes}, \textit{final}, \textit{error}).
La interfaz de usuario se migró de Gradio a una SPA en React~19 con
Vite~8, TypeScript~6, Tailwind~v4 y shadcn/ui. Se incluye un panel
administrativo con recarga en caliente de parámetros LLM y RAG sin
reiniciar el servidor, y una suite de evaluación con métricas nDCG,
Recall@K, Precision@K y MRR sobre un conjunto de referencia (\textit{golden set}).
\end{abstract}

\begin{IEEEkeywords}
agente conversacional, recuperación aumentada por generación, LangGraph,
LlamaIndex, Qdrant, PostgreSQL, FastAPI, SSE, React, \textit{web scraping},
modelos de lenguaje grandes, inteligencia artificial
\end{IEEEkeywords}

% ============================================================
\section{Introducción}
% ============================================================

El Módulo~1 del proyecto estableció una base sólida: un corpus textual de 40+
páginas extraídas de \url{https://valledellili.org/} mediante \textit{web
scraping}, convertidas a Markdown con metadatos YAML y consultables mediante
BM25~\cite{robertson2009bm25}. Si bien ese enfoque demostró la viabilidad del
RAG léxico, presentó limitaciones inherentes: incapacidad para capturar
similitud semántica, recuperación de un único documento sin diversificación y
ausencia de memoria conversacional entre turnos.

El Módulo~2 supera estas limitaciones mediante una arquitectura de agente
completo~\cite{lewis2020rag, chase2022langchain} que incorpora: (i)~recuperación
densa con \textit{embeddings} y similitud coseno en Qdrant~\cite{qdrant2024};
(ii)~diversificación MMR~\cite{carbonell1998mmr} y reranking con
\textit{cross-encoder}; (iii)~memoria conversacional multi-turno en
PostgreSQL~\cite{langchainpostgres}; (iv)~orquestación de flujo con
LangGraph~\cite{langgraph2024}; y (v)~interfaz SPA moderna con React~19
y streaming SSE.

La contribución principal de este trabajo es la demostración de una
arquitectura de agente conversacional educativamente completa, que integra
múltiples herramientas heterogéneas bajo un grafo de orquestación unificado,
con trazabilidad de razonamiento expuesta al usuario final y capacidad de
ajuste administrativo en tiempo de ejecución sin reiniciar el servidor.

% ============================================================
\section{Arquitectura General del Sistema}
\label{sec:arq}
% ============================================================

La Figura~\ref{fig:arq} muestra los componentes del sistema y sus
interacciones. El corpus textual, obtenido mediante \textit{web scraping}
(Módulo~1), fluye hacia Qdrant a través del script de ingesta. En tiempo de
inferencia, cada petición del usuario pasa por FastAPI, que delega al grafo
LangGraph, el cual accede a PostgreSQL (memoria) y Qdrant (RAG) según la
herramienta elegida.

\begin{figure*}[ht]
\centering
\begin{tikzpicture}[
    box/.style={
        rectangle, rounded corners=4pt,
        draw=black!60, fill=blue!7,
        minimum width=3.0cm, minimum height=1.05cm,
        font=\small\sffamily, align=center, inner sep=5pt
    },
    dbase/.style={
        rectangle, rounded corners=4pt,
        draw=black!60, fill=green!10,
        minimum width=2.2cm, minimum height=1.05cm,
        font=\small\sffamily, align=center, inner sep=5pt
    },
    arrow/.style={-Stealth, thick, draw=black!60},
    dasharrow/.style={-Stealth, thick, draw=black!45, dashed},
    lbl/.style={font=\tiny\sffamily}
]
% ---- Fila superior: pipeline de ingesta (offline) ----
\node[box] at (0,    2.6) (raw)
    {\textbf{data/raw/}\\[1pt]\footnotesize HTML + JSON};
\node[box] at (4.0,  2.6) (md)
    {\textbf{data/markdown/}\\[1pt]\footnotesize Markdown + YAML};
\node[box] at (8.2,  2.6) (idx)
    {\textbf{Ingesta}\\[1pt]\footnotesize chunks + embeddings\\[-2pt]\footnotesize upsert idempotente};
\node[dbase] at (13.0, 2.6) (qdrant)
    {\textbf{Qdrant}\\[1pt]\footnotesize v1.18.0};

% ---- Fila inferior: runtime ----
\node[box] at (0,    0) (fe)
    {\textbf{Frontend}\\[1pt]\footnotesize React 19 + Vite 8\\[-2pt]\footnotesize SSE client};
\node[box] at (4.0,  0) (api)
    {\textbf{FastAPI}\\[1pt]\footnotesize Uvicorn\\[-2pt]\footnotesize SSE + Admin};
\node[box] at (8.2,  0) (lg)
    {\textbf{Grafo LangGraph}\\[1pt]\footnotesize 6 nodos secuenciales\\[-2pt]\footnotesize 3 herramientas};
\node[dbase] at (13.0, 0) (pg)
    {\textbf{PostgreSQL}\\[1pt]\footnotesize 17.9-alpine};

% ---- Flechas de ingesta ----
\draw[arrow] (raw) -- node[above,lbl]{scraping} (md);
\draw[arrow] (md)  -- node[above,lbl]{export}   (idx);
\draw[arrow] (idx) -- node[above,lbl]{upsert}   (qdrant);

% ---- Flechas de runtime ----
\draw[arrow]    (fe)  -- node[above,lbl]{POST}  (api);
\draw[dasharrow](fe.south) -- ++(0,-0.30) -| node[below,lbl,pos=0.25]{SSE} ($(api.south)-(0.4,0)$);
\draw[arrow]    (api) -- node[above,lbl]{grafo}  (lg);
\draw[arrow]    (lg)  -- node[above,lbl]{memoria}(pg);

% ---- Flecha LangGraph → Qdrant (vertical entre filas) ----
\draw[arrow] (lg.north) -- node[right,lbl]{rag\_denso} (qdrant.south);

% ---- Etiqueta admin ----
\node[lbl,above=0.08cm of api]{\textit{/api/admin}};

\end{tikzpicture}
\caption{Arquitectura del sistema Módulo~2. El corpus (izquierda) se indexa
\textit{offline} en Qdrant (fila superior). En inferencia (fila inferior),
Frontend $\to$ FastAPI $\to$ LangGraph accede a Qdrant (\code{rag\_denso})
y a PostgreSQL (memoria conversacional).}
\label{fig:arq}
\end{figure*}

La Tabla~\ref{tab:stack} resume el \textit{stack} tecnológico empleado.

\begin{table}[ht]
\centering
\caption{Stack tecnológico del Módulo 2}
\label{tab:stack}
\setlength{\tabcolsep}{3.5pt}
\begin{tabular}{@{}lll@{}}
\toprule
\textbf{Capa} & \textbf{Tecnología} & \textbf{Versión} \\
\midrule
Orquestación   & LangGraph                         & 0.2.x      \\
Tools y memoria& LangChain + langchain-postgres    & 0.3.x      \\
RAG denso      & LlamaIndex + qdrant-client        & 0.12.x     \\
Vector store   & Qdrant                            & v1.18.0    \\
OLTP           & PostgreSQL + SQLAlchemy async     & 17.9 / 2.x \\
Migraciones    & Alembic                           & 1.13.x     \\
Backend HTTP   & FastAPI + Uvicorn + sse-starlette & 0.115.x    \\
Frontend       & React 19 + Vite 8 + TypeScript 6  & —          \\
UI components  & Tailwind v4 + shadcn/ui           & —          \\
Streaming UI   & fetch + ReadableStream + Zod      & —          \\
Contenedores   & Docker Compose                    & —          \\
\bottomrule
\end{tabular}
\end{table}

% ============================================================
\section{Corpus y Proceso de Ingesta}
\label{sec:corpus}
% ============================================================

\subsection{Continuidad con el Módulo 1}

El \textit{web scraping} de \url{https://valledellili.org/} y la conversión a
Markdown con \textit{front matter} YAML descritos en el Módulo~1 permanecen
sin cambios. El directorio \code{data/markdown/valledellili-org/} (40+ archivos)
constituye la fuente de verdad textual del sistema.

\subsection{Ingesta en Qdrant}

El script \code{scripts/indexar\_corpus\_qdrant.py} transforma el corpus textual
en vectores y los carga en Qdrant mediante el proceso de la
Figura~\ref{fig:ingesta}.

\begin{figure}[ht]
\centering
\begin{tikzpicture}[
    node distance=0.38cm,
    box/.style={
        rectangle, rounded corners=3pt,
        draw=black!60, fill=blue!7,
        text width=6.8cm, minimum height=0.72cm,
        font=\small\sffamily, align=center,
        inner sep=4pt
    },
    arrow/.style={-Stealth, thick, draw=black!60}
]
    \node[box] (p1)
        {\textbf{1. Descubrimiento}\quad
         \textit{glob *.md en data/markdown/}};
    \node[box,below=of p1] (p2)
        {\textbf{2. Parseo}\quad
         \textit{front matter YAML + cuerpo Markdown}};
    \node[box,below=of p2] (p3)
        {\textbf{3. Chunking}\quad
         \textit{SentenceSplitter (chunk=1024, overlap=128)}};
    \node[box,below=of p3] (p4)
        {\textbf{4. Enriquecimiento de metadatos}\quad
         \textit{tipo\_pagina, especialidad, sedes, headings}};
    \node[box,below=of p4] (p5)
        {\textbf{5. Embeddings}\quad
         \textit{texto $\rightarrow$ vector $\rightarrow$ normalización L2}};
    \node[box,below=of p5] (p6)
        {\textbf{6. Upsert idempotente en Qdrant}\quad
         \textit{id = UUID5(SHA-256(chunk\_text))}};
    \draw[arrow] (p1)--(p2);
    \draw[arrow] (p2)--(p3);
    \draw[arrow] (p3)--(p4);
    \draw[arrow] (p4)--(p5);
    \draw[arrow] (p5)--(p6);
\end{tikzpicture}
\caption{Pipeline de ingesta: de Markdown a puntos vectoriales en Qdrant.}
\label{fig:ingesta}
\end{figure}

La idempotencia se garantiza mediante el identificador de punto:
\[
\text{id\_chunk} = \mathrm{UUID5}(\mathcal{N},\; \mathrm{SHA256}(\text{chunk\_text}))
\]
donde $\mathcal{N}$ es un \textit{namespace} fijo del proyecto. Un upsert sobre
un chunk sin cambios solo reescribe el vector y el \textit{payload} con valores
idénticos, sin duplicar datos.

Cada punto en Qdrant almacena en su \textit{payload}: \code{archivo},
\code{titulo}, \code{source\_url}, \code{tipo\_pagina}, \code{especialidad},
\code{sedes}, \code{chunk\_index} y un campo \code{texto} con los primeros 500
caracteres del fragmento.

\subsection{Configuración de Embeddings}

El sistema soporta dos proveedores de \textit{embeddings} intercambiables
(Tabla~\ref{tab:emb}), seleccionables mediante la variable
\code{EMBEDDING\_PROVIDER}.

\begin{table}[ht]
\centering
\caption{Proveedores de embeddings soportados}
\label{tab:emb}
\setlength{\tabcolsep}{3.5pt}
\begin{tabular}{@{}llcc@{}}
\toprule
\textbf{Proveedor} & \textbf{Modelo por defecto} & \textbf{Dims} & \textbf{Tipo} \\
\midrule
OpenAI      & text-embedding-3-small & 1536 & API (nube)  \\
HuggingFace & all-MiniLM-L6-v2       & 384  & Local (CPU) \\
\bottomrule
\end{tabular}
\end{table}

% ============================================================
\section{Módulo de Recuperación Densa}
\label{sec:rag}
% ============================================================

\subsection{Similitud Coseno y Búsqueda en Qdrant}

La similitud entre el vector de consulta $\mathbf{q}$ y cada vector de documento
$\mathbf{d}_j$ se calcula como:

\begin{equation}
\mathrm{sim}(\mathbf{q}, \mathbf{d}_j)
= \frac{\mathbf{q} \cdot \mathbf{d}_j}{\|\mathbf{q}\|_2\,\|\mathbf{d}_j\|_2}
\label{eq:coseno}
\end{equation}

Los vectores se normalizan a norma unitaria en la ingesta, por lo que la
búsqueda se reduce a un producto escalar, lo que optimiza el cómputo en Qdrant.
La clase \code{RecuperadorDenso} (en \code{src/rag/recuperador\_denso.py})
solicita inicialmente $K_0 = \max(K_\text{ini}, K)$ candidatos
($K_\text{ini} = 20$ por defecto), que luego se filtran y refinan.

\subsection{Filtro por Umbral y Filtros de Payload}

Tras obtener los candidatos, se descartan los fragmentos cuya similitud no
alcanza el umbral mínimo:
\[
\mathcal{F} = \{(\mathbf{d}_j, s_j) : s_j \geq s_{\min}\}, \quad s_{\min} = 0.25
\]
Adicionalmente, la búsqueda puede restringirse a subconjuntos del corpus
mediante filtros sobre el campo \code{tipo\_pagina} del \textit{payload}
(e.g.\ \code{["institucional"]}, \code{["sede"]}, \code{["ficha\_medico"]}),
infiriendo la intención de la consulta con heurísticas de expresión regular.

\subsection{Reranking con Cross-Encoder (Opcional)}

Cuando \code{RAG\_RERANKER\_HABILITADO=true}, los primeros $N_r$ candidatos
se puntúan con un \textit{cross-encoder} (modelo por defecto:
\code{BAAI/bge-reranker-base}~\cite{bge2023}):

\begin{equation}
\text{score\_rerank}(q, d) = \mathrm{CrossEncoder}(q \oplus d)
\label{eq:rerank}
\end{equation}

donde $\oplus$ denota la concatenación del par consulta-documento. A diferencia
de los \textit{embeddings} bi-encoder, el \textit{cross-encoder} modela la
interacción completa entre consulta y documento en una única pasada, produciendo
puntuaciones más precisas a costa de mayor latencia.

\subsection{Diversificación MMR}

Para evitar fragmentos redundantes en la respuesta, se aplica Maximal Marginal
Relevance~\cite{carbonell1998mmr}. Dado el conjunto de candidatos $\mathcal{F}$
y el subconjunto de documentos ya seleccionados $\mathcal{S}$, el siguiente
fragmento elegido es:

\begin{equation}
d^* = \underset{d \in \mathcal{F} \setminus \mathcal{S}}{\arg\max}
\left[
  \lambda \cdot \mathrm{sim}(d, \mathbf{q})
  - (1-\lambda) \cdot \max_{d' \in \mathcal{S}} \mathrm{sim}(d, d')
\right]
\label{eq:mmr}
\end{equation}

donde $\lambda \in [0,1]$ equilibra la relevancia respecto a la consulta y la
diversidad respecto a los documentos ya seleccionados. El valor por defecto
$\lambda = 0.5$ trata ambos criterios simétricamente.

\subsection{Pipeline Completo de Recuperación}

El orden de aplicación de las etapas de post-proceso es:

\begin{enumerate}
    \item \textbf{Reranker} sobre los primeros $N_r$ candidatos densos (si habilitado).
    \item \textbf{MMR} sobre el pool resultante (si habilitado).
    \item \textbf{Top-$K$ final}: se retornan los $K$ fragmentos con mayor puntuación.
\end{enumerate}

La salida es un objeto \code{SalidaRecuperacionRagDenso} con dos campos clave:
\code{respuesta\_contexto} (texto concatenado de los fragmentos, con cabeceras
de identificación) y \code{fuentes} (lista de \code{FuenteRagDenso} con
\code{score\_denso}, \code{score\_final} y metadatos de trazabilidad).

% ============================================================
\section{Orquestación con LangGraph}
\label{sec:langgraph}
% ============================================================

\subsection{Estado Compartido}

El grafo opera sobre un estado tipado \code{EstadoAgente} (TypedDict de
Python), que actúa como memoria de corto plazo dentro de un turno. Los campos
clave se muestran en la Tabla~\ref{tab:estado}.

\begin{table}[ht]
\centering
\caption{Campos principales del estado del grafo LangGraph}
\label{tab:estado}
\setlength{\tabcolsep}{3pt}
\begin{tabular}{@{}lll@{}}
\toprule
\textbf{Campo} & \textbf{Tipo} & \textbf{Descripción} \\
\midrule
\code{pregunta}            & \code{str}           & Entrada del usuario \\
\code{session\_id}         & \code{str}           & ID de sesión canónico \\
\code{mensajes\_historial} & \code{list[BaseMsg]} & Ventana de memoria \\
\code{intencion}           & \code{str}           & factual / listado / conteo \\
\code{tool\_decidida}      & \code{str | None}    & Herramienta elegida \\
\code{resultado\_tool}     & \code{dict}          & Salida de la herramienta \\
\code{fuentes}             & \code{list[dict]}    & Chunks RAG (si aplica) \\
\code{pensamientos}        & \code{list[dict]}    & Traza del router (SSE) \\
\code{respuesta\_final}    & \code{str}           & Texto del compositor \\
\bottomrule
\end{tabular}
\end{table}

El campo \code{pensamientos} usa el \textit{reducer} \code{operator.add}, que
agrega (en lugar de sobreescribir) los eventos de razonamiento producidos por
cada nodo, permitiendo transmitirlos progresivamente al cliente.

\subsection{Nodos y Flujo del Grafo}

El grafo contiene seis nodos secuenciales sin aristas condicionales
explícitas. Las decisiones de ramificación se toman dentro de cada nodo
mediante lógica condicional interna. La Figura~\ref{fig:grafo} muestra el
flujo completo.

\begin{figure}[ht]
\centering
\begin{tikzpicture}[
    node distance=0.36cm,
    nodo/.style={
        rectangle, rounded corners=3pt,
        draw=black!60, fill=orange!10,
        text width=7.0cm, minimum height=0.68cm,
        font=\small\sffamily, align=center,
        inner sep=4pt
    },
    arrow/.style={-Stealth, thick, draw=black!60}
]
    \node[nodo] (n1)
        {\textbf{N1: cargar\_memoria}\quad
         \textit{Ventana PostgreSQL (días + turnos)}};
    \node[nodo,below=of n1] (n2)
        {\textbf{N2: inferir\_intencion}\quad
         \textit{Regex puro: factual / listado / conteo}};
    \node[nodo,below=of n2] (n3)
        {\textbf{N3: decidir\_tool}\quad
         \textit{Heurística O LLM router + bind\_tools}};
    \node[nodo,below=of n3] (n4)
        {\textbf{N4: ejecutar\_tool}\quad
         \textit{faq\_estructurada / rag\_denso / listar\_estructurado}};
    \node[nodo,below=of n4] (n5)
        {\textbf{N5: componer\_respuesta}\quad
         \textit{LLM compositor con contexto + historial}};
    \node[nodo,below=of n5] (n6)
        {\textbf{N6: persistir\_turno}\quad
         \textit{INSERT HumanMessage + AIMessage en chat\_history}};
    \draw[arrow] (n1)--(n2);
    \draw[arrow] (n2)--(n3);
    \draw[arrow] (n3)--(n4);
    \draw[arrow] (n4)--(n5);
    \draw[arrow] (n5)--(n6);
\end{tikzpicture}
\caption{Grafo LangGraph del agente: seis nodos secuenciales. Las decisiones
de herramienta se toman dentro del nodo N3 y el fallback a RAG desde
N4 si el listado no retorna resultados.}
\label{fig:grafo}
\end{figure}

\subsubsection{Nodo N1 — cargar\_memoria}
Carga la ventana de historial del usuario desde PostgreSQL filtrando por fecha
(\code{HISTORIAL\_DIAS\_MAX}, defecto 7 días) y número de turnos
(\code{HISTORIAL\_TURNOS\_MAX}, defecto 20), usando
\code{PostgresChatMessageHistory} de langchain-postgres.

\subsubsection{Nodo N2 — inferir\_intencion}
Clasifica la consulta sin invocar al LLM mediante expresiones regulares:
\begin{itemize}
    \item \textbf{conteo}: patrones como \textit{``cuántos'', ``cantidad de''}.
    \item \textbf{listado}: patrones como \textit{``lista'', ``enumera''}.
    \item \textbf{factual}: caso por defecto (resto de consultas).
\end{itemize}
Esta clasificación local reduce la latencia en rutas predecibles a menos de 1\,ms.

\subsubsection{Nodo N3 — decidir\_tool}
Si la intención es \textit{listado} o \textit{conteo}, extrae filtros de
\textit{payload} mediante heurísticas (especialidades, sedes, tipo de página)
y genera directamente el \textit{tool call} a \code{listar\_estructurado} sin
consultar al LLM. En caso contrario, invoca al LLM router con
\code{bind\_tools} sobre las definiciones de herramientas del meta-prompt.
El \textit{fallback} predeterminado cuando el router no emite
\textit{tool\_calls} es \code{rag\_denso}.

\subsubsection{Nodo N4 — ejecutar\_tool}
Invoca la herramienta decidida y captura su salida. Si
\code{listar\_estructurado} retorna \code{conteo = 0}, se ejecuta un
\textit{fallback} automático a \code{rag\_denso} con la pregunta original,
garantizando que el usuario reciba siempre una respuesta contextualizada.

\subsubsection{Nodo N5 — componer\_respuesta}
Construye el mensaje al LLM compositor con: instrucciones institucionales,
nombre del usuario (trato personalizado), historial previo y el contexto
JSON serializado de la herramienta. La respuesta se transmite en
\textit{streaming} token a token al cliente.

\subsubsection{Nodo N6 — persistir\_turno}
Almacena el par pregunta-respuesta en \code{chat\_history} (PostgreSQL) con
metadatos de la herramienta utilizada y las fuentes RAG si aplican.

% ============================================================
\section{Las Tres Herramientas del Agente}
\label{sec:tools}
% ============================================================

\subsection{FAQ Estructurada}

La herramienta \code{faq\_estructurada} responde consultas frecuentes sin
necesidad de LLM ni de Qdrant, con latencia inferior a 5\,ms. Cada entrada
del archivo \code{data/structured/faqs.json} contiene: \code{id},
\code{intent}, \code{keywords}, \code{pregunta\_canonica} y \code{respuesta}.

La función de puntuación combina dos señales:

\[
s_\text{kw} = \frac{|\{k \in \text{keywords} : \text{tokens}(k) \subseteq \text{tokens}(q)\}|}{|\text{keywords}|}
\]
\[
s_\text{can} = \frac{|\text{tokens}(p_c) \cap \text{tokens}(q)|}{|\text{tokens}(p_c)|}
\]
\[
s_\text{efec} = \max(s_\text{kw},\; s_\text{can})
\]

donde $q$ es la consulta, $p_c$ es la pregunta canónica y los tokens se
normalizan (NFKD + minúsculas, sin stopwords). Si $s_\text{efec} \geq
\tau_\text{faq}$ (umbral por defecto $= 0.6$), se retorna la respuesta
asociada; en caso contrario el router elige otra herramienta.

Un mecanismo de caché basado en \code{mtime} del archivo JSON permite
actualizar las FAQs sin reiniciar el servidor.

\subsection{RAG Denso}

Implementa el pipeline de recuperación descrito en la Sección~\ref{sec:rag}.
La herramienta acepta los parámetros \code{consulta},
\code{filtros\_tipo\_pagina} (opcionales) y parámetros RAG inyectados desde
el \code{RuntimeAgenteBundle}. El contexto retornado es un bloque de texto con
fragmentos delimitados y sus fuentes.

\subsection{Listar Estructurado}

Realiza un \textit{scroll} sobre el \textit{payload} de Qdrant con filtros
exactos (\code{tipo\_pagina}, \code{especialidad}, \code{sedes}), sin computar
similitud vectorial. Esto permite responder preguntas de agregación como
\textit{``¿cuántos pediatras tienen en Cali?''} con el conteo exacto de puntos
que cumplen los filtros, en lugar de aproximaciones semánticas. La salida
incluye los campos \code{conteo}, \code{items} (lista de registros
deduplicados), \code{muestra\_truncada} y \code{filtros\_aplicados}.

La Tabla~\ref{tab:tools} compara las tres herramientas.

\begin{table}[ht]
\centering
\caption{Comparación de las tres herramientas del agente}
\label{tab:tools}
\setlength{\tabcolsep}{3pt}
\footnotesize
\begin{tabular}{@{}llccp{1.5cm}@{}}
\toprule
\textbf{Herramienta} & \textbf{Backend} & \textbf{LLM} & \textbf{Latencia} & \textbf{Caso de uso} \\
\midrule
\code{faq\_estr.}          & JSON curado   & No & $<$5\,ms    & Preguntas frecuentes \\
\code{rag\_denso}          & Qdrant        & No & 50--300\,ms & Búsqueda semántica   \\
\code{listar\_estr.}       & Qdrant scroll & No & 30--200\,ms & Conteo y listados    \\
\bottomrule
\end{tabular}
\normalsize
\end{table}

% ============================================================
\section{Memoria Conversacional en PostgreSQL}
\label{sec:memoria}
% ============================================================

\subsection{Esquema de Base de Datos}

El sistema mantiene tres tablas en PostgreSQL administradas con SQLAlchemy 2
async y Alembic. La Tabla~\ref{tab:db} resume su propósito.

\begin{table}[ht]
\centering
\caption{Tablas OLTP del sistema}
\label{tab:db}
\setlength{\tabcolsep}{3pt}
\begin{tabular}{@{}lll@{}}
\toprule
\textbf{Tabla} & \textbf{Gestión} & \textbf{Propósito} \\
\midrule
\code{usuarios}        & Alembic   & Registro de usuarios con UUID \\
\code{chat\_history}   & LangChain & Historial de mensajes por sesión \\
\code{config\_admin\_m2} & Alembic & Configuración hot-reload (singleton) \\
\bottomrule
\end{tabular}
\end{table}

La tabla \code{usuarios} almacena \code{id} (UUID), \code{documento\_identidad}
(único), \code{nombre} y marcas de tiempo. La tabla \code{chat\_history} es
creada en \textit{runtime} por \code{PostgresChatMessageHistory.create\_tables}
y almacena cada mensaje como un objeto JSONB con tipo, contenido y metadatos
(herramienta usada, fuentes RAG).

\subsection{Ventana de Memoria y Control Optimista}

La ventana de memoria se configura con dos parámetros independientes:
\code{HISTORIAL\_DIAS\_MAX} (defecto 7 días) y \code{HISTORIAL\_TURNOS\_MAX}
(defecto 20 turnos). Ambas restricciones se aplican en la misma consulta SQL,
garantizando que el contexto de historial no desborde la ventana del LLM.

La tabla \code{config\_admin\_m2} implementa \textit{locking} optimista
mediante un campo \code{version}: el endpoint \code{PATCH /api/admin/config}
solo actualiza si la versión recibida coincide con la almacenada, retornando
HTTP~409 en caso de conflicto.

% ============================================================
\section{API HTTP y Protocolo SSE}
\label{sec:api}
% ============================================================

\subsection{Autenticación Liviana}

El endpoint \code{POST /api/sesiones} recibe \code{documento\_identidad} y
\code{nombre}, crea o recupera el registro en \code{usuarios}, y emite una
cookie HTTP-only con el \code{session\_id} (\code{user:\{uuid\}}). Las
peticiones posteriores se autentican mediante la cookie, la cabecera
\code{X-Session-Id} o el parámetro \code{session\_id} (en ese orden de
precedencia).

\subsection{Streaming SSE Extendido}

El endpoint \code{POST /api/agente/stream} retorna un \code{EventSourceResponse}
con seis tipos de eventos (Tabla~\ref{tab:sse}). El cliente React los procesa
progresivamente: primero muestra el razonamiento del router (\textit{pensamiento}),
luego los metadatos de la herramienta ejecutada (\textit{herramienta}), seguido
de los tokens del compositor en tiempo real (\textit{token}), las fuentes RAG
(\textit{fuentes}, si aplica) y finalmente el evento de cierre (\textit{final}).

\begin{table}[ht]
\centering
\caption{Eventos SSE del agente conversacional}
\label{tab:sse}
\setlength{\tabcolsep}{3pt}
\begin{tabular}{@{}lll@{}}
\toprule
\textbf{Evento} & \textbf{Cuándo se emite} & \textbf{Payload clave} \\
\midrule
\textit{pensamiento} & Tras N3 (decidir\_tool)    & herramienta\_candidata, razón \\
\textit{herramienta} & Tras N4 (ejecutar\_tool)   & nombre, latencia\_ms, resultados \\
\textit{token}       & Por cada token del N5      & motor, texto \\
\textit{fuentes}     & Si RAG retorna chunks       & lista de FuenteRagDenso \\
\textit{final}       & Al finalizar N6             & texto, latencia\_ms, modelo \\
\textit{error}       & Si excepción en el grafo    & código, mensaje \\
\bottomrule
\end{tabular}
\end{table}

Un \textit{watchdog} asíncrono monitorea la desconexión del cliente cada
250\,ms y aborta el streaming si este se desconecta, evitando cómputo
innecesario en el servidor.

\subsection{Panel Administrativo y Hot-Reload}

El router \code{/api/admin} (protegido por la cabecera \code{X-Admin-Key})
expone los endpoints de la Tabla~\ref{tab:admin}. El mecanismo de
\textit{hot-reload} funciona así: el endpoint \code{PATCH /api/admin/config}
actualiza la tabla \code{config\_admin\_m2}; la siguiente petición SSE lee la
configuración actualizada y reconstruye el \code{RuntimeAgenteBundle} sin
reiniciar Uvicorn. Los parámetros ajustables incluyen: modelo LLM del router y
compositor, temperatura, top-p, parámetros RAG (top-k, score mínimo, MMR,
reranker) y los prompts institucionales.

\begin{table}[ht]
\centering
\caption{Endpoints del panel administrativo}
\label{tab:admin}
\setlength{\tabcolsep}{3pt}
\begin{tabular}{@{}lll@{}}
\toprule
\textbf{Método} & \textbf{Ruta} & \textbf{Acción} \\
\midrule
GET   & \code{/api/admin/config}    & Configuración efectiva fusionada \\
PATCH & \code{/api/admin/config}    & Actualización parcial (hot-reload) \\
GET   & \code{/api/admin/usuarios}  & Listado paginado (doc\_id enmascarado) \\
GET   & \code{/api/admin/metricas}  & Agregados: sesiones, total chats \\
\bottomrule
\end{tabular}
\end{table}

% ============================================================
\section{Interfaz de Usuario: React + Vite}
\label{sec:frontend}
% ============================================================

\subsection{Migración desde Gradio}

El Módulo~1 empleaba Gradio para la interfaz de usuario. En el Módulo~2, la
interfaz se migró a una SPA (\textit{Single Page Application}) desarrollada con
React~19 y Vite~8, que ofrece mayor control sobre la experiencia de usuario,
soporte completo de TypeScript~6 estricto, componentes accesibles con
shadcn/ui~\cite{shadcnui} sobre Radix~UI y soporte nativo de modo
oscuro/claro con \code{next-themes}.

\subsection{Cliente SSE Personalizado}

El \code{EventSource} nativo del navegador no soporta peticiones POST con cuerpo
JSON. Para mantener la semántica RESTful del endpoint, se implementó un cliente
SSE propio en \code{frontend/src/lib/sseClient.ts} usando la combinación
\code{fetch} + \code{ReadableStream} + \code{TextDecoderStream}:

\begin{lstlisting}[language=JavaScript,
    caption={Esquema del cliente SSE personalizado},
    label={lst:sse}]
async function* streamSseEvents(url, body, headers) {
  const res = await fetch(url, {
    method: "POST",
    headers: { "Content-Type": "application/json",
               ...headers },
    body: JSON.stringify(body),
    credentials: "include",
  });
  const reader = res.body
    .pipeThrough(new TextDecoderStream())
    .pipeThrough(new EventStreamTransformStream())
    .getReader();
  while (true) {
    const { done, value } = await reader.read();
    if (done) break;
    yield { tipo: value.event,
            payload: JSON.parse(value.data) };
  }
}
\end{lstlisting}

Los eventos se validan con esquemas Zod para garantizar tipado seguro en tiempo
de ejecución. El proxy de Vite en desarrollo reenvía \code{/api/*} al backend
en \code{localhost:8000}, eliminando la necesidad de configurar CORS durante
el desarrollo.

\subsection{Componentes Principales}

\begin{itemize}
    \item \textbf{IdentificationPage}: formulario de documento e identidad antes
          del chat; llama \code{POST /api/sesiones} y almacena el contexto de
          sesión.
    \item \textbf{ChatPage} + \textbf{MessageList}: historial de mensajes con
          renderizado Markdown (react-markdown + remark-gfm) y resaltado de
          sintaxis (shiki). Las fuentes RAG se muestran como tabla colapsable.
    \item \textbf{AdminDashboardPage}: panel de configuración hot-reload con
          selectores de modelo, sliders de temperatura/top-p, toggles de MMR
          y reranker, y editores de prompts.
    \item \textbf{ThemeToggle} + modo oscuro/claro: persiste en
          \code{localStorage} y respeta la preferencia del sistema operativo.
\end{itemize}

% ============================================================
\section{Despliegue con Docker Compose}
\label{sec:deploy}
% ============================================================

El archivo \code{docker-compose.yml} define los servicios de la
Tabla~\ref{tab:docker}. El servicio \code{db-init} ejecuta
\code{alembic upgrade head} como condición previa al arranque de la API,
garantizando que las migraciones se apliquen antes de aceptar tráfico.

\begin{table}[ht]
\centering
\caption{Servicios del stack Docker Compose}
\label{tab:docker}
\setlength{\tabcolsep}{3pt}
\begin{tabular}{@{}llll@{}}
\toprule
\textbf{Servicio} & \textbf{Imagen} & \textbf{Puerto (host)} & \textbf{Rol} \\
\midrule
\code{postgres} & postgres:17.9-alpine   & 15432 & OLTP + memoria \\
\code{qdrant}   & qdrant/qdrant:v1.18.0  & 6333, 6334 & Vector store \\
\code{pgweb}    & sosedoff/pgweb:0.16.2  & 8081  & UI admin Postgres \\
\code{db-init}  & (imagen API)           & ---   & Migraciones Alembic \\
\code{api}      & (multi-stage build)    & 8000  & FastAPI + frontend \\
\bottomrule
\end{tabular}
\end{table}

El \textit{build} multi-stage compila el frontend React y lo sirve como
archivos estáticos desde FastAPI, eliminando la necesidad de un servidor web
adicional en producción. Los datos persisten en volúmenes nombrados
\code{postgres-data} y \code{qdrant-storage}.

% ============================================================
\section{Evaluación}
\label{sec:eval}
% ============================================================

\subsection{Dataset de Referencia}

Para la evaluación del módulo RAG se construyó un \textit{golden set} en
formato JSONL almacenado en \code{data/eval/golden\_set\_rag.jsonl}. Cada
entrada contiene la consulta, el archivo fuente esperado y, opcionalmente,
el fragmento de texto de referencia. Adicionalmente, el archivo
\code{tests/qa/preguntas\_evaluacion.yml} (23 preguntas del Módulo~1) permanece
como referencia de cobertura temática.

\subsection{Métricas de Recuperación}

El script \code{scripts/eval\_metricas\_rag.py} calcula las métricas de la
Tabla~\ref{tab:metricas} para los $K$ fragmentos recuperados.

\begin{table}[ht]
\centering
\caption{Métricas de evaluación del sistema RAG}
\label{tab:metricas}
\setlength{\tabcolsep}{3pt}
\begin{tabular}{@{}lll@{}}
\toprule
\textbf{Métrica} & \textbf{Granularidad} & \textbf{Definición} \\
\midrule
Hit@K      & Chunk    & 1 si algún chunk ∈ top-K coincide con ground truth \\
Precision@K & Doc     & $|\text{docs recuperados} \cap \text{gt}| / K$ \\
Recall@K    & Doc     & $|\text{docs recuperados} \cap \text{gt}| / |\text{gt}|$ \\
MRR         & Chunk   & $1 / \text{rango del primer acierto}$ \\
nDCG@K      & Ranking & $\mathrm{DCG}@K / \mathrm{iDCG}@K$ \\
\bottomrule
\end{tabular}
\end{table}

Donde:
\begin{equation}
\mathrm{DCG}@K = \sum_{i=1}^{K} \frac{\mathrm{rel}_i}{\log_2(i+1)}
\label{eq:dcg}
\end{equation}
\begin{equation}
\mathrm{iDCG}@K = \sum_{i=1}^{\min(K,\,|\mathrm{gt}|)} \frac{1}{\log_2(i+1)}
\label{eq:idcg}
\end{equation}
siendo $\mathrm{rel}_i \in \{0,1\}$ si el $i$-ésimo fragmento recuperado
pertenece al conjunto de referencia.

\subsection{Escenarios de Prueba E2E}

La suite \code{tests/e2e/test\_escenarios\_modulo2.py} cubre cuatro escenarios
con \code{MOCK\_LLM=1} (router determinista) para verificar el comportamiento
del grafo sin dependencia de la API de OpenAI (Tabla~\ref{tab:e2e}).

\begin{table}[ht]
\centering
\caption{Escenarios E2E del Módulo 2}
\label{tab:e2e}
\setlength{\tabcolsep}{3pt}
\footnotesize
\begin{tabular}{@{}llp{2.4cm}@{}}
\toprule
\textbf{Escenario} & \textbf{Herramienta} & \textbf{Verificación} \\
\midrule
A — RAG     & \code{rag\_denso}        & Evento \textit{fuentes} con chunks \\
B — Memoria & \code{rag\_denso} × 2   & Historial persiste entre turnos \\
C — FAQ     & \code{faq\_estructurada} & \code{faq\_match}\\
            &                          & \code{\_encontrado=true} \\
D — Mixto   & FAQ + RAG + FAQ          & Secuencia y estado coherentes \\
\bottomrule
\end{tabular}
\normalsize
\end{table}

\subsection{Resultados}

La Tabla~\ref{tab:resultados} presentará los resultados cuantitativos de
recuperación y generación una vez ejecutado el protocolo sobre el corpus
definitivo.

\begin{table}[ht]
\centering
\caption{Resultados de evaluación RAG por configuración}
\label{tab:resultados}
\setlength{\tabcolsep}{3pt}
\begin{tabular}{@{}lcccc@{}}
\toprule
\textbf{Configuración} & \textbf{Hit@5} & \textbf{Recall@5}
                       & \textbf{nDCG@5} & \textbf{Latencia (ms)} \\
\midrule
Solo denso (base)       & --- & --- & --- & --- \\
Denso + MMR             & --- & --- & --- & --- \\
Denso + Reranker        & --- & --- & --- & --- \\
Denso + Reranker + MMR  & --- & --- & --- & --- \\
\bottomrule
\end{tabular}
\end{table}

\textit{Nota: los valores se completarán con los resultados del \textit{golden set}
definitivo ejecutado sobre el corpus de producción.}

% ============================================================
\section{Discusión}
\label{sec:disc}
% ============================================================

\subsection{Comparación con el Módulo 1 (BM25)}

La migración de BM25 a recuperación densa representa un cambio cualitativo
en la capacidad de comprensión semántica del sistema. BM25 requería coincidencia
léxica exacta; el sistema de Módulo~2 puede recuperar documentos relevantes ante
consultas con sinónimos, abreviaciones o paráfrasis. Sin embargo, este avance
tiene un costo: la recuperación densa depende de un proveedor de
\textit{embeddings} (en la configuración por defecto, la API de OpenAI), lo que
introduce una dependencia externa y un costo por token. La opción de HuggingFace
local mitiga esta dependencia a expensas de mayor latencia de ingesta.

\subsection{Ventajas del Diseño con Tres Herramientas}

La segmentación en \code{faq\_estructurada}, \code{rag\_denso} y
\code{listar\_estructurado} responde a una observación empírica: distintos
tipos de consultas tienen perfiles de información óptimos. Las preguntas
frecuentes (horarios, contacto, datos institucionales) son mejor respondidas
por respuestas curadas y deterministas. Las preguntas conceptuales o clínicas
se benefician de la recuperación semántica. Los listados y conteos requieren
exactitud sobre el conjunto de datos completo, no aproximaciones.

\subsection{Limitaciones}

\begin{enumerate}
    \item \textbf{Sincronización corpus-vectores}: cambios en
          \code{data/markdown/} requieren reindexación explícita para reflejarse
          en Qdrant.
    \item \textbf{Latencia del reranker en CPU}: el modelo
          \code{BAAI/bge-reranker-base} puede tardar 300--800\,ms por turno en
          CPU, lo que lo hace inadecuado en producción sin GPU.
    \item \textbf{Costo de API}: la inferencia del router y los embeddings
          dependen de OpenAI en la configuración por defecto.
    \item \textbf{Retención de datos}: el historial conversacional contiene
          datos personales; la ventana temporal configurable mitiga la
          acumulación indefinida pero no sustituye una política de privacidad
          formal.
\end{enumerate}

% ============================================================
\section{Conclusiones}
\label{sec:conc}
% ============================================================

Este trabajo presentó un agente conversacional completo para la Fundación
Valle del Lili que supera las limitaciones del sistema BM25 del Módulo~1
mediante seis contribuciones técnicas:

\begin{enumerate}
    \item \textbf{RAG denso con pipeline configurable}: similitud coseno en
          Qdrant, filtro por umbral, reranking con \textit{cross-encoder} y
          diversificación MMR, todo controlable desde variables de entorno o
          el panel administrativo en tiempo de ejecución.
    \item \textbf{Orquestación con LangGraph}: grafo de seis nodos secuenciales
          con clasificación de intención mediante regex (sin LLM), selección
          de herramienta por heurísticas o \textit{tool-calling}, y mecanismo
          de \textit{fallback} encadenado.
    \item \textbf{Tres herramientas complementarias}: FAQ determinista, RAG
          semántico y listado estructurado por \textit{payload}, cubriendo
          los principales patrones de consulta institucional.
    \item \textbf{Memoria persistente}: ventana de historial en PostgreSQL
          con control temporal y de turnos, integrada de forma transparente
          en el estado del grafo.
    \item \textbf{API SSE con seis tipos de eventos}: trazabilidad completa del
          razonamiento del agente expuesta al cliente, desde la decisión de
          herramienta hasta las fuentes RAG.
    \item \textbf{Panel administrativo con hot-reload}: ajuste de parámetros
          LLM, RAG y prompts sin reiniciar el servidor, con control de
          versión optimista para entornos con múltiples operadores.
\end{enumerate}

La arquitectura demuestra que es posible construir un sistema de asistencia
conversacional institucional robusto, trazable y ajustable en producción,
combinando componentes de código abierto (LangGraph, LlamaIndex, Qdrant,
PostgreSQL, React) con modelos de lenguaje de acceso comercial.

% ============================================================
\begin{thebibliography}{00}
% ============================================================

\bibitem{lewis2020rag}
P.~Lewis, E.~Perez, A.~Piktus \textit{et al.},
``Retrieval-Augmented Generation for Knowledge-Intensive NLP Tasks,''
en \textit{Advances in Neural Information Processing Systems (NeurIPS)},
vol.~33, pp.~9459--9474, 2020.

\bibitem{robertson2009bm25}
S.~Robertson y H.~Zaragoza,
``The Probabilistic Relevance Framework: BM25 and Beyond,''
\textit{Foundations and Trends in Information Retrieval},
vol.~3, n.º~4, pp.~333--389, 2009.

\bibitem{chase2022langchain}
H.~Chase,
``LangChain: Building applications with LLMs through composability,''
\textit{GitHub repository}, 2022.
[En línea]. Disponible en: \url{https://github.com/langchain-ai/langchain}

\bibitem{langgraph2024}
LangChain AI,
``LangGraph: Build stateful, multi-actor applications with LLMs,''
2024.
[En línea]. Disponible en: \url{https://github.com/langchain-ai/langgraph}

\bibitem{qdrant2024}
Qdrant,
``Qdrant: Vector Database for the next generation of AI applications,''
2024.
[En línea]. Disponible en: \url{https://qdrant.tech}

\bibitem{carbonell1998mmr}
J.~Carbonell y J.~Goldstein,
``The Use of MMR, Diversity-Based Reranking for Reordering Documents
and Producing Summaries,''
en \textit{Proc.\ 21st ACM SIGIR}, pp.~335--336, 1998.

\bibitem{bge2023}
S.~Xiao, Z.~Liu, P.~Zhang, N.~Muennighoff,
``C-Pack: Packaged Resources To Advance General Chinese Embedding,''
\textit{arXiv preprint arXiv:2309.07597}, 2023.

\bibitem{langchainpostgres}
LangChain AI,
``langchain-postgres: PostgreSQL integration for LangChain,''
2024.
[En línea]. Disponible en: \url{https://github.com/langchain-ai/langchain-postgres}

\bibitem{llamaindex2024}
J.~Liu,
``LlamaIndex: Data Framework for LLM Applications,''
\textit{GitHub repository}, 2022.
[En línea]. Disponible en: \url{https://github.com/run-llama/llama_index}

\bibitem{openai2024gpt4o}
OpenAI,
``GPT-4o System Card,''
2024.
[En línea]. Disponible en: \url{https://openai.com/index/gpt-4o-system-card}

\bibitem{fastapi2024}
S.~Ramírez,
``FastAPI: Modern, fast (high-performance) web framework for building APIs,''
2024.
[En línea]. Disponible en: \url{https://fastapi.tiangolo.com}

\bibitem{shadcnui}
shadcn,
``shadcn/ui: Beautifully designed components built with Radix UI and Tailwind,''
2024.
[En línea]. Disponible en: \url{https://ui.shadcn.com}

\end{thebibliography}

\end{document}
